\uuid{xu6O}
\exo7id{1845}
\titre{exo7 1845}
\auteur{ridde}
\organisation{exo7}
\datecreate{1999-11-01}
\isIndication{false}
\isCorrection{false}
\chapitre{Equation aux dérivées partielles}
\sousChapitre{Equation aux dérivées partielles}
\cpge{mpsi-mp2i.fonctions-deux-variables.b, pcsi.fonctions-deux-variables.b, ptsi.fonctions-deux-variables.b, bcpst1.derivees-primitives-integrales.b, bcpst1.fonctions-deux-variables.c, mp-mpi.calcul-differentiel-optimisation.a}
\module{Equation différentielle}
\niveau{L3}
\difficulte{}

\contenu{
\question{
R\'esoudre \`a l'aide des coordonn\'ees polaires l'\'equation aux d\'eriv\'ees partielles :
$$x\frac{\partial f}{\partial x} (x, y) + y\frac{\partial f}{\partial y} (x, y) =
\sqrt{x^{2} + y^{2}} $$
}
}
